% GENERATED FILE. Edit canonical lessons, then run npm run generate:books.
% canonical-source-hash: fnv1a64:e6fdafeb8b1fdafb
% canonical-lessons: IT-C365-giustizia, IT-C365-ingiusto, IT-C365-ingiustizia, IT-C365-pregiudizio, IT-C365-stereotipo

\chapter{Giusto e ingiusto}
\label{ch:it-giusto-e-ingiusto}
\hlchaptermodality{365}

\begin{chapteropening}
\textbf{By the end of this chapter:} \emph{I can talk about justice and injustice, prejudice and stereotypes.}

Complete the last lesson of chapter 365: I can talk about justice and injustice, prejudice and stereotypes.
\end{chapteropening}

\section[la giustizia]{la giustizia — the justice}

\label{lesson:IT-C365-giustizia}



\begin{quote}
Before the new one: say the Italian for to demonstrate, then the Italian for the banner.
\end{quote}

\subsection*{You'll want to know: la giustizia}
\textbf{la giustizia} — \textquotedblleft{}the justice\textquotedblright{}.

\begin{quote}
Le vittime chiedono giustizia.
\end{quote}

\begin{quote}
\emph{The victims are asking for justice.}
\end{quote}

\emph{Fare giustizia} — to see justice done.

\begin{cousinweb}
\emph{giustizia} is a learned borrowing from Latin \emph{iūstitia} \textquotedblleft{}justice, fairness\textquotedblright{} (Lewis \& Short), from \emph{iūstus} \textquotedblleft{}fair, right\textquotedblright{}.
\end{cousinweb}

\begin{grammarlens}[title={What you've built}]
One more word for telling, explaining and describing at B1.
\end{grammarlens}

\subsection*{Your turn}
\begin{itemize}
\raggedright
  \item \emph{Say it:} \emph{la giustizia}
  \item \emph{Say it:} \emph{la giustizia}, once more
  \item \emph{Say it:} a sentence of your own with \emph{la giustizia}
  \item \emph{Recall:} say the Italian for to demonstrate, then the Italian for the banner, then say \emph{la giustizia} again
\end{itemize}

\begin{tcolorbox}[breakable,colback=teal!4,colframe=teal!35!black,title={Before you move on}]
What does la giustizia mean? (\textquotedblleft{}The justice\textquotedblright{}.) Say it once more.
\end{tcolorbox}
\section[ingiusto]{ingiusto — unfair, unjust}

\label{lesson:IT-C365-ingiusto}



\begin{quote}
Before the new one: say the Italian for the banner, then the Italian for the justice.
\end{quote}

\subsection*{You'll want to know: ingiusto}
\textbf{ingiusto} — \textquotedblleft{}unfair, unjust\textquotedblright{}.

\begin{quote}
È ingiusto dare la colpa a tutta la classe per un errore.
\end{quote}

\begin{quote}
\emph{It is unfair to blame the whole class for one mistake.}
\end{quote}

\emph{In-} \textquotedblleft{}not\textquotedblright{} + \emph{giusto}.

\begin{cousinweb}
\emph{ingiusto} comes from Latin \emph{iniūstus} \textquotedblleft{}unjust\textquotedblright{} (Lewis \& Short), built from \emph{in-} \textquotedblleft{}not\textquotedblright{} and \emph{iūstus} \textquotedblleft{}fair, right\textquotedblright{}.
\end{cousinweb}

\begin{grammarlens}[title={What you've built}]
One more word for telling, explaining and describing at B1.
\end{grammarlens}

\subsection*{Your turn}
\begin{itemize}
\raggedright
  \item \emph{Say it:} \emph{ingiusto}
  \item \emph{Say it:} \emph{ingiusto}, once more
  \item \emph{Say it:} a sentence of your own with \emph{ingiusto}
  \item \emph{Recall:} say the Italian for the banner, then the Italian for the justice, then say \emph{ingiusto} again
\end{itemize}

\begin{tcolorbox}[breakable,colback=teal!4,colframe=teal!35!black,title={Before you move on}]
What does ingiusto mean? (\textquotedblleft{}Unfair, unjust\textquotedblright{}.) Say it once more.
\end{tcolorbox}
\section[l'ingiustizia]{l'ingiustizia — the injustice}

\label{lesson:IT-C365-ingiustizia}



\begin{quote}
Before the new one: say the Italian for the justice, then the Italian for unfair.
\end{quote}

\subsection*{You'll want to know: l'ingiustizia}
\textbf{l'ingiustizia} — \textquotedblleft{}the injustice\textquotedblright{}.

\begin{quote}
Questa è una grande ingiustizia.
\end{quote}

\begin{quote}
\emph{This is a great injustice.}
\end{quote}

\emph{Subire un'ingiustizia} — to suffer an injustice.

\begin{cousinweb}
\emph{ingiustizia} is a learned borrowing from Latin \emph{iniūstitia} \textquotedblleft{}injustice\textquotedblright{} (Lewis \& Short), from \emph{iniūstus}, built on \emph{iūstus} \textquotedblleft{}fair, right\textquotedblright{}.
\end{cousinweb}

\begin{grammarlens}[title={What you've built}]
One more word for telling, explaining and describing at B1.
\end{grammarlens}

\subsection*{Your turn}
\begin{itemize}
\raggedright
  \item \emph{Say it:} \emph{l'ingiustizia}
  \item \emph{Say it:} \emph{l'ingiustizia}, once more
  \item \emph{Say it:} a sentence of your own with \emph{l'ingiustizia}
  \item \emph{Recall:} say the Italian for the justice, then the Italian for unfair, then say \emph{l'ingiustizia} again
\end{itemize}

\begin{tcolorbox}[breakable,colback=teal!4,colframe=teal!35!black,title={Before you move on}]
What does l'ingiustizia mean? (\textquotedblleft{}The injustice\textquotedblright{}.) Say it once more.
\end{tcolorbox}
\section[il pregiudizio]{il pregiudizio — the prejudice}

\label{lesson:IT-C365-pregiudizio}



\begin{quote}
Before the new one: say the Italian for unfair, then the Italian for the injustice.
\end{quote}

\subsection*{You'll want to know: il pregiudizio}
\textbf{il pregiudizio} — \textquotedblleft{}the prejudice\textquotedblright{}.

\begin{quote}
Bisogna conoscere le persone senza pregiudizi.
\end{quote}

\begin{quote}
\emph{You have to get to know people without prejudice.}
\end{quote}

Literally a judgement made before.

\begin{cousinweb}
\emph{pregiudizio} is a learned borrowing from Latin \emph{praeiūdicium} \textquotedblleft{}a judgement made beforehand\textquotedblright{} (Lewis \& Short), built from \emph{prae} \textquotedblleft{}before\textquotedblright{} and \emph{iūdicium} \textquotedblleft{}judgement\textquotedblright{}, on \emph{iūdex} \textquotedblleft{}judge\textquotedblright{}.
\end{cousinweb}

\begin{grammarlens}[title={What you've built}]
One more word for telling, explaining and describing at B1.
\end{grammarlens}

\subsection*{Your turn}
\begin{itemize}
\raggedright
  \item \emph{Say it:} \emph{il pregiudizio}
  \item \emph{Say it:} \emph{il pregiudizio}, once more
  \item \emph{Say it:} a sentence of your own with \emph{il pregiudizio}
  \item \emph{Recall:} say the Italian for unfair, then the Italian for the injustice, then say \emph{il pregiudizio} again
\end{itemize}

\begin{tcolorbox}[breakable,colback=teal!4,colframe=teal!35!black,title={Before you move on}]
What does il pregiudizio mean? (\textquotedblleft{}The prejudice\textquotedblright{}.) Say it once more.
\end{tcolorbox}
\section[lo stereotipo]{lo stereotipo — the stereotype}

\label{lesson:IT-C365-stereotipo}



\begin{quote}
Before the new one: say the Italian for the injustice, then the Italian for the prejudice.
\end{quote}

\subsection*{You'll want to know: lo stereotipo}
\textbf{lo stereotipo} — \textquotedblleft{}the stereotype\textquotedblright{}.

\begin{quote}
Gli stereotipi sugli italiani sono tanti: pizza, mandolino...
\end{quote}

\begin{quote}
\emph{There are many stereotypes about Italians: pizza, mandolin...}
\end{quote}

Stress: \emph{stereòtipo}.

\begin{cousinweb}
\emph{stereotipo} is usually traced to French \emph{stéréotype}, a printing word built from Greek \emph{stereós} \textquotedblleft{}solid\textquotedblright{} and \emph{týpos} \textquotedblleft{}a blow, an impression\textquotedblright{}, the word behind Latin \emph{typus} \textquotedblleft{}a figure, an image\textquotedblright{} (Lewis \& Short): a fixed plate that prints the same image again and again.
\end{cousinweb}

\begin{grammarlens}[title={What you've built}]
One more word for telling, explaining and describing at B1.
\end{grammarlens}

\subsection*{Your turn}
\begin{itemize}
\raggedright
  \item \emph{Say it:} \emph{lo stereotipo}
  \item \emph{Say it:} \emph{lo stereotipo}, once more
  \item \emph{Say it:} a sentence of your own with \emph{lo stereotipo}
  \item \emph{Recall:} say the Italian for the injustice, then the Italian for the prejudice, then say \emph{lo stereotipo} again
\end{itemize}

\begin{tcolorbox}[breakable,colback=teal!4,colframe=teal!35!black,title={Before you move on}]
What does lo stereotipo mean? (\textquotedblleft{}The stereotype\textquotedblright{}.) Say it once more.
\end{tcolorbox}
